\section{Texto}
\hspace{0.55cm}Para resaltar textos se usan los siguientes comandos:
\begin{itemize}
    \item \textbf{Negrita}: se usa el comando \textbf{\textbackslash{textbf\{\}}} para hacer el texto negrita.
    \item \underline{Subrayado}: se usa el comando \textbf{\textbackslash{underline\{\}}} para hacer una subrayar el texto.
    \item \textit{Cursiva}: se usa el comando \textbf{\textbackslash{textit\{\}}} para hacer el texto cursiva.
\end{itemize}
\begin{lstlisting}
    \begin{document}
        \textbf{Este texto es negrita.}
        \textit{Este texto es cursiva.}
        \underline{Este texto está subrayado.}
    \end{document}
\end{lstlisting}



\section{Comentarios}
\hspace{0.55cm}Para poner un comentario se usa el símbolo \%:
\begin{center}
    \textit{\% Esto es un comentario.}
\end{center}
%Esto es un comentario.

Si lo que se busca es comentar varias líneas de código al mismo tiempo sin ir una a una, podemos utilizar el paquete \textbf{comment} y utilizar su comando \textbf{comment}:
\begin{lstlisting}
    \usepackage{verbatim}   % Paquete para comentar varias líneas a la vez.
    
    \begin{document}
        \begin{comment}
        esto
        es
        un
        bloque
        de
        texto
        comentado
        \end{comment}
        
        hola mundo
    \end{document}
\end{lstlisting}
\begin{comment}
    Hola mamá esto no se verá.
\end{comment}

Si buscamos el camino más rápido, Overleaf (la plataforma online donde actualmente estamos escribiendo este documento) proporciona el \textit{shorcut} \textbf{Ctrl + /} para comentar varias líneas de texto ya seleccionadas con el cursor. Este \textit{shorcut} puede variar dependiendo el IDE o editor de texto donde se esté desarrollando el proyecto \LaTeX.



\section{Escritura de código en \LaTeX}
\hspace{0.55cm}Algo muy interesante a la hora de trabajar en este ambiente, es el lograr escribir código de programación que se vea elegante y tabulado para su mejor comprensión, esto lo logramos con el paquete \textbf{listing}.
\begin{center}
    \textit{\textbackslash{usepackage\{listings\}}}
\end{center}
    
Lo más difícil (o sencillo), es personalizar o configurar este paquete para que se mire como nosotros deseemos, para ello utilizamos el comando \textbf{lstset\{\}}, este debe ser posicionado antes del comienzo del documento, y dentro de sus llaves podemos insertar muchas entradas o atributos que configurarán el resultado final del paquete. Algo que quiero mencionar es que este paquete soporta automáticamente la escritura y estilo de varios lenguajes de programación (Python, Java, C, C++, HTML, CSS), esto es un tipo de entrada o atributo que podemos configurar en \textit{lstset}, dejo a continuación la configuración usada en este documento y otra que podría utilizarse para escribir código de CSharp:

\textbf{Caso 1, configuración completa de este archivo}:
\begin{lstlisting}
    % Personalización de la fuente para el código.
    \lstset{
        basicstyle=\ttfamily,		% Utiliza la fuente tttfamily, en especial el paquete inconsolata.
        frame=none,					% Quita el marco al cuadro flotante que contiene el código o texto.
        columns=fullflexible,		% Ajusta el cuadro flotante al tamaño del texto del documento.
        breaklines=true,				% Ajusta el texto dentro del contenedor.
        inputencoding = utf8,		% Admite caracteres del código UTF8.
        extendedchars = true,		% Soporte para caracteres especiales.
        showstringspaces = false,	% Quita los guiones bajos predeterminados de los espacios en cadenas de texto.
        escapebegin = \obeyspaces,	% Complemento de la entrada anterior.
        literate =					% Acepta los siguientes caracteres especiales fuera de UTF8.
            {á}{{\'a}}1 {é}{{\'e}}1 {í}{{\'i}}1 {ó}{{\'o}}1 {ú}{{\'u}}1
            {Á}{{\'A}}1 {É}{{\'E}}1 {Í}{{\'I}}1 {Ó}{{\'O}}1 {Ú}{{\'U}}1
            {ñ}{{\~n}}1 {Ñ}{{\~N}}1,
    }
\end{lstlisting}

Puedes agregar el comando \textit{\textbackslash{footnotesize}} para hacer un poco más pequeña la letra del código, esto va justo enseguida de \textit{\textbackslash{ttfamily}} en la entrada \textit{basicstyle}.

\textbf{Caso 2, configuración rápida con lenguaje}:
\begin{lstlisting}
    \lstdefinestyle{sharpc}{language=[Sharp]C, frame=lr, rulecolor=\color{blue!80!black}}
    
    \lstset{style=sharpc}
\end{lstlisting}

Para comenzar a trabajarlo utilizamos los comandos \textbf{\textbackslash{begin\{\}}} y \textbf{\textbackslash{end\{\}}}, dentro de sus llaves escribimos la palabra \textbf{lstlisting}: \\
\textit{
    \textbackslash{begin}\{lstlisting\} \\
    Esto es código!!! \\
    \textbackslash{end}\{lstlisting\}
}